\documentclass[10pt,a4paper]{amsart}
\usepackage{amsmath,amssymb,amsthm,graphicx}
% For dealing with references we use the comment environment
\usepackage{verbatim}
\newenvironment{reference}{\comment}{\endcomment}
%\newenvironment{reference}{}{}
\newenvironment{slogan}{\comment}{\endcomment}
\newenvironment{history}{\comment}{\endcomment}
% For commutative diagrams we use Xy-pic
\usepackage[all]{xy}

% We use 2cell for 2-commutative diagrams.
\xyoption{2cell}
\UseAllTwocells
\theoremstyle{plain}
\newtheorem{theorem}[subsection]{Theorem}
\newtheorem{proposition}[subsection]{Proposition}
\newtheorem{lemma}[subsection]{Lemma}

\theoremstyle{definition}
\newtheorem{definition}[subsection]{Definition}
\newtheorem{example}[subsection]{Example}
\newtheorem{exercise}[subsection]{Exercise}
\newtheorem{situation}[subsection]{Situation}

\theoremstyle{remark}
\newtheorem{remark}[subsection]{Remark}
\newtheorem{remarks}[subsection]{Remarks}

\numberwithin{equation}{subsection}
% Macros
%
\def\lim{\mathop{\mathrm{lim}}\nolimits}
\def\colim{\mathop{\mathrm{colim}}\nolimits}
\def\Spec{\mathop{\mathrm{Spec}}}
\def\Hom{\mathop{\mathrm{Hom}}\nolimits}
\def\Ext{\mathop{\mathrm{Ext}}\nolimits}
\def\SheafHom{\mathop{\mathcal{H}\!\mathit{om}}\nolimits}
\def\SheafExt{\mathop{\mathcal{E}\!\mathit{xt}}\nolimits}
\def\Sch{\mathit{Sch}}
\def\Mor{\mathop{\mathrm{Mor}}\nolimits}
\def\Ob{\mathop{\mathrm{Ob}}\nolimits}
\def\Sh{\mathop{\mathit{Sh}}\nolimits}
\def\NL{\mathop{N\!L}\nolimits}
\def\CH{\mathop{\mathrm{CH}}\nolimits}
\def\proetale{{pro\text{-}\acute{e}tale}}
\def\etale{{\acute{e}tale}}
\def\QCoh{\mathit{QCoh}}
\def\Ker{\mathop{\mathrm{Ker}}}
\def\Im{\mathop{\mathrm{Im}}}
\def\Coker{\mathop{\mathrm{Coker}}}
\def\Coim{\mathop{\mathrm{Coim}}}

% Boxtimes
%
\DeclareMathSymbol{\boxtimes}{\mathbin}{AMSa}{"02}

%
% Macros for moduli stacks/spaces
%
\def\QCohstack{\mathcal{QC}\!\mathit{oh}}
\def\Cohstack{\mathcal{C}\!\mathit{oh}}
\def\Spacesstack{\mathcal{S}\!\mathit{paces}}
\def\Quotfunctor{\mathrm{Quot}}
\def\Hilbfunctor{\mathrm{Hilb}}
\def\Curvesstack{\mathcal{C}\!\mathit{urves}}
\def\Polarizedstack{\mathcal{P}\!\mathit{olarized}}
\def\Complexesstack{\mathcal{C}\!\mathit{omplexes}}
% \Pic is the operator that assigns to X its picard group, usage \Pic(X)
% \Picardstack_{X/B} denotes the Picard stack of X over B
% \Picardfunctor_{X/B} denotes the Picard functor of X over B
\def\Pic{\mathop{\mathrm{Pic}}\nolimits}
\def\Picardstack{\mathcal{P}\!\mathit{ic}}
\def\Picardfunctor{\mathrm{Pic}}
\def\Deformationcategory{\mathcal{D}\!\mathit{ef}}
\usepackage{cleveref}
\crefname{originallemma}{original lemma}{original lemmas}
\Crefname{originallemma}{Original lemma}{Original lemmas}
\makeatletter
\newlabel{lemma-change-big-site}{{03FP}{source}{}{sourcepointer-1}{}}
\newlabel{lemma-change-big-site@cref}{{[originallemma][0][]03FP}{[1][0][]source}{}{}{}}
\makeatother
\crefname{originalequation}{original equation}{original equations}
\Crefname{originalequation}{Original equation}{Original equations}
\makeatletter
\newlabel{sets-equation-bound}{{046U}{source}{}{sourcepointer-2}{}}
\newlabel{sets-equation-bound@cref}{{[originalequation][0][]046U}{[1][0][]source}{}{}{}}
\makeatother
\crefname{originallemma}{original lemma}{original lemmas}
\Crefname{originallemma}{Original lemma}{Original lemmas}
\makeatletter
\newlabel{sites-lemma-bigger-site}{{00XU}{source}{}{sourcepointer-3}{}}
\newlabel{sites-lemma-bigger-site@cref}{{[originallemma][0][]00XU}{[1][0][]source}{}{}{}}
\makeatother
\crefname{originalsection}{original section}{original sections}
\Crefname{originalsection}{Original section}{Original sections}
\makeatletter
\newlabel{topologies-section-change-alpha}{{022I}{source}{}{sourcepointer-4}{}}
\newlabel{topologies-section-change-alpha@cref}{{[originalsection][0][]022I}{[1][0][]source}{}{}{}}
\makeatother
\crefname{originalsection}{original section}{original sections}
\Crefname{originalsection}{Original section}{Original sections}
\makeatletter
\newlabel{sets-section-categories-schemes}{{000H}{source}{}{sourcepointer-5}{}}
\newlabel{sets-section-categories-schemes@cref}{{[originalsection][0][]000H}{[1][0][]source}{}{}{}}
\makeatother
\crefname{originallemma}{original lemma}{original lemmas}
\Crefname{originallemma}{Original lemma}{Original lemmas}
\makeatletter
\newlabel{sets-lemma-bound-finite-type}{{04T7}{source}{}{sourcepointer-6}{}}
\newlabel{sets-lemma-bound-finite-type@cref}{{[originallemma][0][]04T7}{[1][0][]source}{}{}{}}
\makeatother
\crefname{originallemma}{original lemma}{original lemmas}
\Crefname{originallemma}{Original lemma}{Original lemmas}
\makeatletter
\newlabel{sets-lemma-bound-by-covering}{{0AHK}{source}{}{sourcepointer-7}{}}
\newlabel{sets-lemma-bound-by-covering@cref}{{[originallemma][0][]0AHK}{[1][0][]source}{}{}{}}
\makeatother
\crefname{originallemma}{original lemma}{original lemmas}
\Crefname{originallemma}{Original lemma}{Original lemmas}
\makeatletter
\newlabel{sets-lemma-bound-size}{{000Q}{source}{}{sourcepointer-8}{}}
\newlabel{sets-lemma-bound-size@cref}{{[originallemma][0][]000Q}{[1][0][]source}{}{}{}}
\makeatother
\crefname{originallemma}{original lemma}{original lemmas}
\Crefname{originallemma}{Original lemma}{Original lemmas}
\makeatletter
\newlabel{sets-lemma-bound-fppf-covering}{{0AHL}{source}{}{sourcepointer-9}{}}
\newlabel{sets-lemma-bound-fppf-covering@cref}{{[originallemma][0][]0AHL}{[1][0][]source}{}{}{}}
\makeatother
\crefname{originallemma}{original lemma}{original lemmas}
\Crefname{originallemma}{Original lemma}{Original lemmas}
\makeatletter
\newlabel{lemma-space-presentation}{{0262}{source}{}{sourcepointer-10}{}}
\newlabel{lemma-space-presentation@cref}{{[originallemma][0][]0262}{[1][0][]source}{}{}{}}
\makeatother
\crefname{originaltheorem}{original theorem}{original theorems}
\Crefname{originaltheorem}{Original theorem}{Original theorems}
\makeatletter
\newlabel{theorem-presentation}{{02WW}{source}{}{sourcepointer-11}{}}
\newlabel{theorem-presentation@cref}{{[originaltheorem][0][]02WW}{[1][0][]source}{}{}{}}
\makeatother
\title{Lemma fully faithful in spaces.tex}
\author{The Stacks Project Authors}
\date{stacks-project:a04446e57ec1fbc252a871afcec7752fb2807b14; GNU Free Documentation License 1.2 (as published in the Stacks Project COPYING)}
\begin{document}
\maketitle
\noindent\textit{Editorial note.} All mathematical statement and proof text is reproduced verbatim from the author's chapter spaces.tex of the Stacks Project, fixed at commit a04446e57ec1fbc252a871afcec7752fb2807b14. The pinned preamble is supplied as its own file; the wrapper class, the xr-hyper wiring and the \textbackslash{}externaldocument list are omitted because the assembler supplies the class.\par\medskip
\noindent\textit{Editorial note.} Numbering inside this standalone excerpt is editorial, not the full Stacks book print numbering.\par\medskip
\noindent\textit{Editorial note.} The candidate is from the corrected band-I follow-on screen (line C) of the extractor's stacks-candidates line: its author label is absent from the live master in every form searched and it belongs to none of the earlier judged pools. The H2 judgement and its reason were made in this module by reading the author's own statement and complete proof in the fixed checkout. Original candidate key: stacks:spaces.tex:lemma-fully-faithful.\par\medskip
\noindent\textit{Editorial note.} Reference labels in the proof resolve to navigation pointers to the original author source; these pointers do not reproduce proofs of the prerequisites.\par\medskip
\noindent\textit{Allowed standard prerequisites (editorial).}\begin{itemize}
\item Prerequisites named by the author inside the counted span. The counted span is reproduced verbatim, including the author's own citations. Its prerequisites are the earlier results of the same fixed Stacks Project version that the author cites there and in the chapter's lead-in.
\end{itemize}\medskip
\ifdefined\hypertarget\hypertarget{sourcepointer-1}{}\fi
\noindent\textit{Source reference.} Original lemma 03FP: sources/spaces.tex, source lines 2558--2558. Navigation pointer used in the counted span; the span contains the pointer only and the cited result is not reproduced inside it. Registered as a non-load-bearing navigation stub in the established form for Stacks material. No bibliography citation is involved.\par\medskip
\ifdefined\hypertarget\hypertarget{sourcepointer-2}{}\fi
\noindent\textit{Source reference.} Original equation 046U: sources/sets.tex, source lines 346--346. Navigation pointer used in the counted span; the span contains the pointer only and the cited result is not reproduced inside it. Registered as a non-load-bearing navigation stub in the established form for Stacks material. No bibliography citation is involved.\par\medskip
\ifdefined\hypertarget\hypertarget{sourcepointer-3}{}\fi
\noindent\textit{Source reference.} Original lemma 00XU: sources/sites.tex, source lines 4730--4730. Navigation pointer used in the counted span; the span contains the pointer only and the cited result is not reproduced inside it. Registered as a non-load-bearing navigation stub in the established form for Stacks material. No bibliography citation is involved.\par\medskip
\ifdefined\hypertarget\hypertarget{sourcepointer-4}{}\fi
\noindent\textit{Source reference.} Original section 022I: sources/topologies.tex, source lines 3986--3986. Navigation pointer used in the counted span; the span contains the pointer only and the cited result is not reproduced inside it. Registered as a non-load-bearing navigation stub in the established form for Stacks material. No bibliography citation is involved.\par\medskip
\ifdefined\hypertarget\hypertarget{sourcepointer-5}{}\fi
\noindent\textit{Source reference.} Original section 000H: sources/sets.tex, source lines 311--311. Navigation pointer used in the counted span; the span contains the pointer only and the cited result is not reproduced inside it. Registered as a non-load-bearing navigation stub in the established form for Stacks material. No bibliography citation is involved.\par\medskip
\ifdefined\hypertarget\hypertarget{sourcepointer-6}{}\fi
\noindent\textit{Source reference.} Original lemma 04T7: sources/sets.tex, source lines 610--610. Navigation pointer used in the counted span; the span contains the pointer only and the cited result is not reproduced inside it. Registered as a non-load-bearing navigation stub in the established form for Stacks material. No bibliography citation is involved.\par\medskip
\ifdefined\hypertarget\hypertarget{sourcepointer-7}{}\fi
\noindent\textit{Source reference.} Original lemma 0AHK: sources/sets.tex, source lines 794--794. Navigation pointer used in the counted span; the span contains the pointer only and the cited result is not reproduced inside it. Registered as a non-load-bearing navigation stub in the established form for Stacks material. No bibliography citation is involved.\par\medskip
\ifdefined\hypertarget\hypertarget{sourcepointer-8}{}\fi
\noindent\textit{Source reference.} Original lemma 000Q: sources/sets.tex, source lines 525--525. Navigation pointer used in the counted span; the span contains the pointer only and the cited result is not reproduced inside it. Registered as a non-load-bearing navigation stub in the established form for Stacks material. No bibliography citation is involved.\par\medskip
\ifdefined\hypertarget\hypertarget{sourcepointer-9}{}\fi
\noindent\textit{Source reference.} Original lemma 0AHL: sources/sets.tex, source lines 833--833. Navigation pointer used in the counted span; the span contains the pointer only and the cited result is not reproduced inside it. Registered as a non-load-bearing navigation stub in the established form for Stacks material. No bibliography citation is involved.\par\medskip
\ifdefined\hypertarget\hypertarget{sourcepointer-10}{}\fi
\noindent\textit{Source reference.} Original lemma 0262: sources/spaces.tex, source lines 1171--1171. Navigation pointer used in the counted span; the span contains the pointer only and the cited result is not reproduced inside it. Registered as a non-load-bearing navigation stub in the established form for Stacks material. No bibliography citation is involved.\par\medskip
\ifdefined\hypertarget\hypertarget{sourcepointer-11}{}\fi
\noindent\textit{Source reference.} Original theorem 02WW: sources/spaces.tex, source lines 1603--1603. Navigation pointer used in the counted span; the span contains the pointer only and the cited result is not reproduced inside it. Registered as a non-load-bearing navigation stub in the established form for Stacks material. No bibliography citation is involved.\par\medskip
\begin{lemma}
\label{lemma-fully-faithful}
Suppose $\Sch_{fppf}$ is contained in $\Sch'_{fppf}$.
Let $S$ be an object of $\Sch_{fppf}$. Denote
$\textit{Spaces}/S$ the category of algebraic spaces over $S$
defined using $\Sch_{fppf}$. Similarly, denote
$\textit{Spaces}'/S$ the category of algebraic spaces over $S$
defined using $\Sch'_{fppf}$. The construction of
Lemma \ref{lemma-change-big-site}
defines a fully faithful functor
$$
\textit{Spaces}/S \longrightarrow \textit{Spaces}'/S
$$
whose essential image consists of those $X' \in \Ob(\textit{Spaces}'/S)$
such that there exist $U, R \in \Ob((\Sch/S)_{fppf})$\footnote{Requiring the
existence of $R$ is necessary because of our choice of the function $Bound$ in
Sets, Equation (\ref{sets-equation-bound}). The size of the fibre product
$U \times_{X'} U$ can grow faster than $Bound$ in terms of the size of $U$. We
can illustrate this by setting $S = \Spec(A)$, $U = \Spec(A[x_i, i \in I])$ and
$R = \coprod_{(\lambda_i) \in A^I} \Spec(A[x_i, y_i]/(x_i - \lambda_i y_i))$.
In this case the size of $R$ grows like $\kappa^\kappa$ where $\kappa$ is the
size of $U$.} and morphisms
$$
U \longrightarrow X'
\quad\text{and}\quad
R \longrightarrow U \times_{X'} U
$$
in $\Sh((\Sch'/S)_{fppf})$ which are surjective as maps of sheaves
(for example if the displayed morphisms are surjective and \'etale).
\end{lemma}

\begin{proof}
In Sites, Lemma \ref{sites-lemma-bigger-site} we have seen that the functor
$f^{-1} : \Sh((\Sch/S)_{fppf}) \to \Sh((\Sch'/S)_{fppf})$
is fully faithful (see discussion in
Topologies, Section \ref{topologies-section-change-alpha}).
Hence we see that the displayed functor of the lemma is fully faithful.

\medskip\noindent
Suppose that $X' \in \Ob(\textit{Spaces}'/S)$ such that there exists
$U \in \Ob((\Sch/S)_{fppf})$ and a map $U \to X'$ in
$\Sh((\Sch'/S)_{fppf})$ which is surjective as a map of sheaves.
Let $U' \to X'$ be a surjective \'etale morphism with
$U' \in \Ob((\Sch'/S)_{fppf})$. Let $\kappa = \text{size}(U)$, see
Sets, Section \ref{sets-section-categories-schemes}.
Then $U$ has an affine open covering $U = \bigcup_{i \in I} U_i$
with $|I| \leq \kappa$. Observe that $U' \times_{X'} U \to U$ is \'etale
and surjective. For each $i$ we can pick a quasi-compact
open $U'_i \subset U'$ such that $U'_i \times_{X'} U_i \to U_i$
is surjective (because the scheme $U' \times_{X'} U_i$ is the
union of the Zariski opens $W \times_{X'} U_i$ for $W \subset U'$
affine and because $U' \times_{X'} U_i \to U_i$ is \'etale hence open).
Then $\coprod_{i \in I} U'_i \to X$ is surjective \'etale
because of our assumption that $U \to X$ and hence $\coprod U_i \to X$
is a surjection of sheaves (details omitted).
Because $U'_i \times_{X'} U \to U'_i$ is a surjection of sheaves
and because $U'_i$ is quasi-compact,
we can find a quasi-compact open $W_i \subset U'_i \times_{X'} U$
such that $W_i \to U'_i$ is surjective as a map of sheaves
(details omitted). Then $W_i \to U$ is \'etale and we conclude
that $\text{size}(W_i) \leq \text{size}(U)$, see
Sets, Lemma \ref{sets-lemma-bound-finite-type}. By
Sets, Lemma \ref{sets-lemma-bound-by-covering} we conclude
that $\text{size}(U'_i) \leq \text{size}(U)$.
Hence $\coprod_{i \in I} U'_i$ is isomorphic to an object of
$(\Sch/S)_{fppf}$ by Sets, Lemma \ref{sets-lemma-bound-size}.

\medskip\noindent
Now let $X'$, $U \to X'$ and $R \to U \times_{X'} U$ be as in the
statement of the lemma. In the previous paragraph we have seen that
we can find $U' \in \Ob((\Sch/S)_{fppf})$ and a surjective \'etale morphism
$U' \to X'$ in $\Sh((\Sch'/S)_{fppf})$. Then
$U' \times_{X'} U \to U'$ is a surjection of sheaves, i.e., we can find an fppf
covering $\{U'_i \to U'\}$ such that $U'_i \to U'$ factors through
$U' \times_{X'} U \to U'$.
By Sets, Lemma \ref{sets-lemma-bound-fppf-covering}
we can find
$\tilde U \to U'$
which is surjective, flat, and locally of finite presentation,
with $\text{size}(\tilde U) \leq \text{size}(U')$,
such that $\tilde U \to U'$ factors through
$U' \times_{X'} U \to U'$. Then we consider
$$
\xymatrix{
U' \times_{X'} U' \ar[d] &
\tilde U \times_{X'} \tilde U  \ar[l] \ar[d] \ar[r] &
U \times_{X'} U \ar[d] \\
U' \times_S U' & \tilde U \times_S \tilde U \ar[l] \ar[r] & U \times_S U
}
$$
The squares are cartesian. We know the objects of the bottom row
are represented by objects of $(\Sch/S)_{fppf}$. By the result of the
argument of the previous paragraph, the same is true for $U \times_{X'} U$
(as we have the surjection of sheaves $R \to U \times_{X'} U$
by assumption). Since $(\Sch/S)_{fppf}$ is closed under fibre
products (by construction), we see that $\tilde U \times_{X'} \tilde U$
is represented by an object of $(\Sch/S)_{fppf}$. Finally, the
map $\tilde U \times_{X'} \tilde U \to U' \times_{X'} U'$ is
a surjection of fppf sheaves as $\tilde U \to U'$ is so.
Thus we can once more apply the result of the previous paragraph
to conclude that $R' = U' \times_{X'} U'$ is represented by an object
of $(\Sch/S)_{fppf}$. At this point
Lemma \ref{lemma-space-presentation} and
Theorem \ref{theorem-presentation} imply that $X = h_{U'}/h_{R'}$
is an object of $\textit{Spaces}/S$ such that $f^{-1}X \cong X'$
as desired.
\end{proof}

\end{document}
